%%% INTRODUCTION

Clinical and lesion studies suggest that declarative memories initially depend on the hippocampus, but are later transferred to other brain areas \cite{Dudai2015, Squire2015, Sekeres2017}. Some forms of memory eventually become independent of the hippocampus and depend only on a stable representation in the neocortex \cite{Dudai2015, Squire2015, Sekeres2017}. Similarly, procedural memories are consolidated within cortico-striatal networks \cite{Karni1994, Brashers-Krug1996, Dudai2015}.
This process of memory transformation --- termed systems memory consolidation --- is thought to prevent newly acquired memories from overwriting old ones, thereby extending memory retention times (``plasticity-stability dilemma''; \cite{Grossberg1987, Abraham2005, Fusi2005, Leibold2008, Roxin2013}), and to enable a simultaneous acquisition of episodic memories and semantic knowledge of the world \cite{McClelland1995, Kumaran2016}.
While specific neuronal activity patterns, including for example an accelerated replay of recent experiences \cite{Lee2002, Skaggs1996}, are involved in the transfer of memories from hippocampus to neocortex \cite{Diekelmann2010}, the mechanisms underlying systems memory consolidation are not well understood. Specifically, it is unclear how this consolidation-related transfer is shaped by the anatomical structure and the plasticity of the underlying neural circuits.
This poses a substantial obstacle for understanding into which regions memories are consolidated; why some memories are consolidated more rapidly than others \cite{Nadel1997, Tse2007, Brodt2018}; why some memories stay hippocampus dependent, and why and how the character of memories changes over time \cite{Dudai2015}; and whether the consolidation of declarative and non-declarative memories \cite{Karni1994, Brashers-Krug1996, Dudai2015} are two sides of the same coin. These questions are hard to approach within phenomenological theories of systems consolidation such as the standard consolidation theory \cite{Squire1995, McClelland1995}, the multiple trace theory \cite{Nadel1997}, and the trace transformation theory \cite{Winocur2010, Winocur2011}.
Here, we propose a novel mechanistic foundation of the consolidation process that accounts for several experimental observations and that could contribute to understanding the transfer of memories and the reorganisation of memories over time on a neuronal level.

Our focus lies on simple forms of memory that can be phrased as cue-response associations. We assume that such associations are stored in synaptic pathways between an input area --- neurally representing the cue --- and an output area --- neurally representing the response.

Thus, our work relates to feedforward, hetero-associative memory (and is therefore applicable to both declarative and non-declarative memories) rather than recurrent, auto-associative memory (see, e.g., \cite{Hopfield1982, Zenke2015, Tome2020}).

Our central hypothesis --- the parallel pathway theory (PPT) --- is that systems memory consolidation arises naturally from the interplay of two abundantly found neuronal features: parallel synaptic pathways and Hebbian plasticity \cite{VanEssen1992, Malenka2004}.
First, we illustrate the underlying mechanism in a simple hippocampal circuit motif and show that Hebbian plasticity can consolidate previously stored associations into parallel pathways.
Next, we outline the PPT in a mathematical framework for the simplest possible (linear) case. Then we show in simulations that the proposed mechanism is robust to various neuronal nonlinearities; further, the mechanism reproduces the results of a hippocampal lesion study in rodents \cite{Remondes2004}; iterated in a cascade, it can achieve a full consolidation into neocortex and result in power-law forgetting of memories as is observed in psychophysical studies in humans \cite{Wixted2004}.
